
\sourcenote{Genre: core rule (tier 1). The rule's semantic content is
honestly small --- the behavioural pullback is two lines --- and the
section says so rather than inflating it; the substance of any
application is the sequent-level discharge of the compatibility
squares, plus the finite-kernel instance where the transport is proven
against actual chain machinery. The gaps beyond those two strata are
marked loudly.}

\textbf{Setting and claim.}\enspace A simulation witness
$\sigma : S \to S'$ is a per-mode polynomial map whose
\emph{compatibility squares} --- commutation of $\sigma$ with the two
systems' dynamics declarations, mode by mode --- are checked as
inherited sequent leaves. The rule: from
$\mathsf{QSatI}(S', \mathit{spec}')$ and the witness, conclude
$\mathsf{QSatI}(S, \sigma^{*}\mathit{spec}')$ --- certify the
abstraction once, reuse under refinement \cite{lynchvaandrager1995,
capuccimyers2026}.

\begin{bthm}[behavioral pullback]\label{thm:qst-pullback}
Let $\mathrm{Beh}$, $\mathrm{Beh}'$ be the compiled behaviour classes
of $S$, $S'$, and suppose $\sigma$ maps behaviours to behaviours
($\sigma[\mathrm{Beh}] \subseteq \mathrm{Beh}'$ --- the semantic
content of the compatibility squares). If every behaviour of $S'$
satisfies $\mathit{spec}'$, then every behaviour of $S$ satisfies
$\sigma^{-1}[\mathit{spec}']$.
\end{bthm}
\begin{proof}
$\omega \in \mathrm{Beh}$ gives
$\sigma(\omega) \in \mathrm{Beh}' \subseteq \mathit{spec}'$, i.e.\
$\omega \in \sigma^{-1}[\mathit{spec}']$
(\texttt{pullback}). Two lines is the honest size: the rule's
\emph{semantic} content is inclusion-chasing, and everything
substantive in an application is establishing
$\sigma[\mathrm{Beh}] \subseteq \mathrm{Beh}'$ from the checked
squares --- per-mode polynomial identities, discharged by the base
calculus. The hypothesis is not droppable to
``$\sigma$ is any map'': without the squares the image of a concrete
behaviour need not be an abstract behaviour, and the pullback claims
anything at all.
\end{proof}

\textbf{The finite-kernel instance.}\enspace At the finite Markov
tier the witness has teeth the behavioural form leaves abstract: a
\emph{functional simulation} is per-state marginal identity ---
pushing the concrete kernel row at $s$ along $\sigma$ yields the
abstract kernel row at $\sigma(s)$ (a pushforward-kernel witness) ---
and the transport theorems are proven against the chain machinery.

\begin{blemma}[observable pullback]\label{lem:qst-expect}
For any law $\mu$ on the concrete states and observable $g$ on the
abstract states,
$\sum_b (\sigma_{*}\mu)(b) \, g(b) = \sum_s \mu(s) \, g(\sigma(s))$;
and along a pushforward-kernel witness, the time-$n$ expectation of a
pulled-back observable on the concrete chain \emph{equals} the
abstract chain's expectation of $g$ at time $n$ (started from the
pushed initial law). Hence any certified expectation bound proven
against the abstract chain holds for the concrete chain at the
pulled-back observable.
\end{blemma}
\begin{proof}
The first identity is summation over fibres of $\sigma$
(\texttt{expect\_push}). The second composes it with the
marginal-process lemma --- the per-state marginal identities
propagate, by induction on time, to equality of pushed laws at every
time --- giving \texttt{sim\_expect}; the bound transports by
rewriting (\texttt{sim\_expect\_bound}). Note the direction: the
producer proves one bound on the small chain and every refinement
inherits it; nothing is re-proven per refinement beyond the witness
identities themselves.
\end{proof}

\begin{blemma}[region pullback]\label{lem:qst-region}
Along a pushforward-kernel witness with entrywise-nonnegative kernel
and initial law: if the abstract law at time $n$ puts no mass outside
$R'$, the concrete law at time $n$ puts no mass outside
$\sigma^{-1}[R']$.
\end{blemma}
\begin{proof}
For $s$ with $\sigma(s) \notin R'$: the concrete mass at $s$ is at
most the pushed mass at $\sigma(s)$ (one nonnegative term of a
fibre sum), which equals the abstract mass at $\sigma(s)$ (the
marginal-process lemma), which is zero; nonnegativity closes the
squeeze (\texttt{sim\_region}). The nonnegativity hypotheses are
essential, not hygiene: with signed ``mass'' the fibre sum could
cancel --- a negative sibling hiding positive mass at $s$ outside the
preimage region --- and the single-term bound fails. Nonnegativity is
what makes ``no mass outside $R'$'' transfer pointwise through a
many-to-one map.
\end{proof}

\begin{bremark}[refinement is one-sided; a simulation is not a
coupling]\label{rem:qst-onesided}
The witness machinery genuinely shared with the relational
(two-chain) tier is exactly the pushforward-kernel identity and the
marginal-process lemma --- a functional simulation \emph{is} such a
witness. Nothing else transfers: the relational tier's certificates
are joint-kernel data over a product space with absorption and drift
hypotheses, and none of that is simulation transport. A functional
simulation induces a degenerate coupling, but the theorems above
never route through it --- they are pushforward plus summation
identities, full stop. Conflating the two would import hypotheses
(coupling emptiness, absorption) that this rule neither needs nor
checks.
\end{bremark}

\begin{bremark}[the gaps, loud]\label{rem:qst-gaps}
The conditional mark scopes the warrant to the two strata proven
above: behavioural, and finite-kernel. Not yet mechanized: the general
measurable-kernel tier (the finite arguments are summation
identities; their measure-kernel analogues are standard but unwritten
here), and the categorical strengthening under which the composition
rule could cite the lens-composition theorem of
\cite{capuccimyers2026} as a lemma \emph{template} --- the direct
composition proofs never pull back along strict weakenings, so
nothing in this book depends on that strengthening; it is named so
that its absence is a recorded choice, not an oversight. No
end-to-end certificate exercises the rule in this edition. Related
reading, not dependency: an exchangeability quotient is a simulation
map, so aggregate-space reduction can be \emph{read} as an instance
of this rule; it stands on its own proof and cites nothing here.
\end{bremark}

\begin{anchors}
Theorem \ref{thm:qst-pullback} & \texttt{pullback}\\
Lemma \ref{lem:qst-expect} & \texttt{expect\_push}, \texttt{sim\_expect}, \texttt{sim\_expect\_bound}; marginal-process lemma \texttt{push\_law} (chain substrate module)\\
Lemma \ref{lem:qst-region} & \texttt{sim\_region}; law nonnegativity \texttt{law\_nonneg} (chain substrate module)\\
witness form & \texttt{PushKernel} (chain substrate module) --- the per-state marginal identity\\
compatibility squares (cited, not re-proved) & inherited sequent leaves, per mode (main package)\\
\end{anchors}
